% META: {"id": "1SPE-VARALEA-CR-000", "chapitre": "1SPE-VARIABLES-ALEATOIRES", "type_objet": "cours", "sous_type": "ouverture", "titre": "Variables aleatoires", "temps_gabarit": 1, "parcours_duree": {"P1": "14h", "P2": "11h", "P3": "9h"}, "prerequis": ["R1", "R2", "R4", "R5"], "capacites_codes": ["C1", "C2", "C3", "C4", "C5", "C6", "C7"], "mode_creation": "ex_nihilo", "sources_inspiration": [], "status": "generated"}

\section*{Contrat du chapitre}

\subsection*{Ce que tu vas savoir faire}

À la fin de ce chapitre, tu seras capable de :

\begin{itemize}
  \item[\textbf{C1}] Interpréter les événements liés à une variable aléatoire et déterminer sa loi de probabilité.
  \item[\textbf{C2}] Calculer l'espérance $E(X)$, la variance $V(X)$ et l'écart type $\sigma(X)$ d'une variable aléatoire.
  \item[\textbf{C3}] Pour deux à quatre épreuves de Bernoulli indépendantes et identiques, construire un arbre et étudier le nombre de succès.
  \item[\textbf{C4}] Utiliser la linéarité de l'espérance pour calculer $E(aX+b)$.
  \item[\textbf{C5}] Résoudre des problèmes contextualisés (jeux, décisions, assurances) à l'aide de variables aléatoires.
  \item[\textbf{C6}] Simuler une variable aléatoire et des échantillons, puis lire, comprendre et écrire une fonction Python renvoyant la moyenne d'un échantillon de taille $n$.
  \item[\textbf{C7}] Étudier expérimentalement la fluctuation des moyennes d'échantillons autour de l'espérance, et interpréter l'écart observé au regard de $\dfrac{2\sigma}{\sqrt{n}}$.
\end{itemize}

\subsection*{Ce dont tu as besoin avant de commencer}

Ce chapitre s'appuie sur quatre familles de prérequis. Le diagnostic de la page suivante te dira lesquels tu maîtrises déjà.

\begin{itemize}
  \item[\textbf{R1}] \textbf{Probabilités de base} — expérience aléatoire, univers, événements, probabilité d'un événement.
  \item[\textbf{R2}] \textbf{Probabilités conditionnelles} — formule des probabilités totales, indépendance.
  \item[\textbf{R4}] \textbf{Suites et sommes} — notation $\sum$, suites arithmétiques et géométriques.
  \item[\textbf{R5}] \textbf{Puissances} — règles de calcul, puissances entières.
\end{itemize}

Si une fiche R te concerne, lis-la avant d'aborder le cours : chaque notion manquante sera un obstacle au moment où tu en auras le plus besoin.

\subsection*{Temps de travail estimé}

\begin{itemize}
  \item Parcours~\parcoursUn{} (Consolidation) : environ \textbf{14~h}.
  \item Parcours~\parcoursDeux{} (Maîtrise) : environ \textbf{11~h}.
  \item Parcours~\parcoursTrois{} (Approfondissement) : environ \textbf{9~h}.
\end{itemize}

\subsection*{Situation d'accroche — L'assurance smartphone}

Un assureur propose un contrat couvrant la casse d'un smartphone. La prime annuelle est de $50$~euros. En cas de casse, l'assureur rembourse $300$~euros. On estime que la probabilité de casse dans l'année est $\dfrac{1}{10}$.

\begin{itemize}
  \item Si l'on note $X$ le gain de l'assuré (remboursement moins prime), quelles sont les valeurs possibles de $X$ ?
  \item Quelle est la probabilité de chaque valeur ?
  \item En moyenne, sur un grand nombre d'assurés, ce contrat est-il avantageux pour l'assuré ? Pour l'assureur ?
\end{itemize}

On a $X = 250$ (cas de casse : $300 - 50$) avec probabilité $\frac{1}{10}$ et $X = -50$ (pas de casse : $-50$) avec probabilité $\frac{9}{10}$. L'espérance $E(X) = 250 \times \frac{1}{10} + (-50) \times \frac{9}{10} = 25 - 45 = -20$. En moyenne, l'assuré perd $20$~euros par an : le contrat est avantageux pour l'assureur.

Cette situation sera entièrement traitée au Temps~7 (TD fil rouge), une fois les capacités C1 à C5 acquises.

% BEGIN-VERIFY
% from sympy import Rational
%
% # L'accroche : prime 50 EUR, remboursement 300 EUR, casse avec probabilite 1/10.
% prime, remboursement, p_casse = 50, 300, Rational(1, 10)
%
% # Gain de l'ASSURE : il paie toujours la prime, et touche le remboursement
% # seulement en cas de casse.
% loi = {remboursement - prime: p_casse, -prime: 1 - p_casse}
% assert set(loi) == {250, -50}
% assert sum(loi.values()) == 1
%
% esperance_assure = sum(valeur*proba for valeur, proba in loi.items())
% assert esperance_assure == Rational(250, 10) - Rational(50*9, 10)
% assert esperance_assure == -20
%
% # Le contrat est donc DEFAVORABLE a l'assure en moyenne, et favorable a
% # l'assureur : les deux esperances sont opposees.
% esperance_assureur = -esperance_assure
% assert esperance_assureur == 20 > 0
% assert esperance_assure < 0
%
% # Le point d'equilibre : la prime qui annulerait l'esperance de l'assure.
% equitable = remboursement * p_casse
% assert equitable == 30
% loi_equitable = {remboursement - equitable: p_casse, -equitable: 1 - p_casse}
% assert sum(v*p for v, p in loi_equitable.items()) == 0
% END-VERIFY
